\chapter{1993年京湘鄂云琼贵卷（理）}
\section{单选题}

\begin{problem}
函数 \(f(x)=\sin x+\cos x\) 的最小正周期是（\quad）

\choices
    {\(2\pi\)}
    {\(2\sqrt{2}\pi\)}
    {\(\pi\)}
    {\(\frac{\pi}{4}\)}
\end{problem}

\begin{answer}
A
\end{answer}

\begin{solution}
利用两角和公式，
\[
f(x)=\sin x+\cos x=\sqrt{2}\sin\left(x+\frac{\pi}{4}\right)
\]
正弦函数的最小正周期为 \(2\pi\)，平移和伸缩不改变周期，因此 \(f(x)\) 的最小正周期为 \(2\pi\)，故选 A.
\end{solution}

\begin{problem}
如果双曲线的焦距为 \(6\)，两条准线间的距离为 \(4\)，那么该双曲线的离心率为（\quad）

\choices
    {\(\frac{3}{2}\)}
    {\(\frac{\sqrt{3}}{2}\)}
    {\(\frac{\sqrt{6}}{2}\)}
    {\(2\)}
\end{problem}

\begin{answer}
C
\end{answer}

\begin{solution}
设双曲线的实半轴长为 \(a\)，半焦距为 \(c\). 由焦距为 \(6\)，得 \(2c=6\)，即 \(c=3\). 双曲线两条准线的方程为 \(x=\pm\frac{a}{e}=\pm\frac{a^2}{c}\)，所以两条准线间的距离为
\[
2\frac{a^2}{c}=4
\]
代入 \(c=3\) 得 \(a^2=6\)，从而离心率
\[
e=\frac{c}{a}=\frac{3}{\sqrt{6}}=\frac{\sqrt{6}}{2}
\]
故选 C.
\end{solution}

\begin{problem}
和直线 \(3x-4y+5=0\) 关于 \(x\) 轴对称的直线的方程为（\quad）

\choices
    {\(3x+4y-5=0\)}
    {\(3x+4y+5=0\)}
    {\(-3x+4y-5=0\)}
    {\(-3x+4y+5=0\)}
\end{problem}

\begin{answer}
B
\end{answer}

\begin{solution}
关于 \(x\) 轴对称的变换为 \((x,y)\mapsto(x,-y)\). 将原直线方程中的 \(y\) 换成 \(-y\)，得
\[
3x-4(-y)+5=0
\]
即 \(3x+4y+5=0\)，故选 B.
\end{solution}

\begin{problem}
极坐标方程 \(\rho=\frac{4}{3-5\cos\theta}\) 所表示的曲线是（\quad）

\choices
    {焦点到准线距离为 \(\frac{4}{5}\) 的椭圆}
    {焦点到准线距离为 \(\frac{4}{5}\) 的双曲线右支}
    {焦点到准线距离为 \(\frac{4}{3}\) 的椭圆}
    {焦点到准线距离为 \(\frac{4}{3}\) 的双曲线右支}
\end{problem}

\begin{answer}
B
\end{answer}

\begin{solution}
极坐标圆锥曲线的标准形式为
\[
\rho=\frac{ed}{1-e\cos\theta}
\]
其中 \(e\) 是离心率，\(d\) 是焦点到准线的距离. 将题给方程改写为
\[
\rho=\frac{4/3}{1-\frac{5}{3}\cos\theta}
\]
因此 \(e=\frac{5}{3}>1\)，曲线是双曲线；且
\(ed=\frac{4}{3}\)，\(\Longrightarrow\)，\(d=\frac{4/3}{5/3}=\frac45\).
当 \(\theta=0\) 时分母为负，按该极坐标形式对应双曲线的右支，故选 B.
\end{solution}

\begin{problem}
\(y=x^{\frac{3}{5}}\) 在 \([-1,1]\) 上是（\quad）

\choices
    {增函数且是奇函数}
    {增函数且是偶函数}
    {减函数且是奇函数}
    {减函数且是偶函数}
\end{problem}

\begin{answer}
A
\end{answer}

\begin{solution}
因为 \(5\) 是奇数，\(x^{3/5}\) 对每个实数 \(x\) 都有意义. 若 \(x_1<x_2\)，则
\(\sqrt[5]{x_1}<\sqrt[5]{x_2}\)，\(\Longrightarrow\)，\(\left(\sqrt[5]{x_1}\right)^3<\left(\sqrt[5]{x_2}\right)^3\)
所以 \(y=x^{3/5}\) 是增函数. 又
\[
(-x)^{3/5}=-x^{3/5}
\]
所以它是奇函数，故选 A.
\end{solution}

\begin{problem}
\(\displaystyle\lim\limits_{n\to\infty}\frac{5n^2-1}{2n^2-n+5}\) 的值为（\quad）

\choices
    {\(-\frac{1}{5}\)}
    {\(-\frac{5}{2}\)}
    {\(\frac{1}{5}\)}
    {\(\frac{5}{2}\)}
\end{problem}

\begin{answer}
D
\end{answer}

\begin{solution}
分子、分母同除以 \(n^2\)，得
\[
\lim_{n\to\infty}\frac{5n^2-1}{2n^2-n+5}
=\lim_{n\to\infty}\frac{5-\frac1{n^2}}{2-\frac1n+\frac5{n^2}}
=\frac52
\]
故选 D.
\end{solution}

\begin{problem}
集合 \(M=\left\{x\mid x=\frac{k\pi}{2}+\frac{\pi}{4},\ k\in\Z\right\}\)，\(N=\left\{x\mid x=\frac{k\pi}{4}+\frac{\pi}{2},\ k\in\Z\right\}\)，则（\quad）

\choices
    {\(M=N\)}
    {\(M\supseteq N\)}
    {\(M\subseteq N\)}
    {\(M\cap N=\varnothing\)}
\end{problem}

\begin{answer}
C
\end{answer}

\begin{solution}
将两集合中的元素分别写成以 \(\frac\pi4\) 为单位的形式：
\(M=\left\{\frac{(2k+1)\pi}{4}\mid k\in\Z\right\}\)，\(N=\left\{\frac{(k+2)\pi}{4}\mid k\in\Z\right\} =\left\{\frac{m\pi}{4}\mid m\in\Z\right\}\).
因此 \(M\subseteq N\). 又 \(0\in N\) 而 \(0\notin M\)，所以两集合不相等，故选 C.
\end{solution}

\begin{problem}
\(\sin 20^\circ\cos 70^\circ+\sin 10^\circ\sin 50^\circ\) 的值是（\quad）

\choices
    {\(\frac{1}{4}\)}
    {\(\frac{\sqrt{3}}{2}\)}
    {\(\frac{1}{2}\)}
    {\(\frac{\sqrt{3}}{4}\)}
\end{problem}

\begin{answer}
A
\end{answer}

\begin{solution}
由积化和差公式，
\[
\sin10^\circ\sin50^\circ
=\frac12\left(\cos40^\circ-\cos60^\circ\right)
\]
又 \(\cos70^\circ=\sin20^\circ\)，所以
\[
\sin20^\circ\cos70^\circ+\sin10^\circ\sin50^\circ
=\frac{1-\cos40^\circ}{2}
+\frac{\cos40^\circ-\frac12}{2}
=\frac14
\]
故选 A.
\end{solution}

\begin{problem}
参数方程 \(x=\left|\cos\frac{\theta}{2}+\sin\frac{\theta}{2}\right|\)，\(y=\frac{1}{2}(1+\sin\theta)\)（\(0<\theta<2\pi\)）表示（\quad）

\choices
    {双曲线的一支，这支过点 \(\left(1,\frac{1}{2}\right)\)}
    {抛物线的一部分，这部分过 \(\left(1,\frac{1}{2}\right)\)}
    {双曲线的一支，这支过点 \(\left(-1,\frac{1}{2}\right)\)}
    {抛物线的一部分，这部分过 \(\left(-1,\frac{1}{2}\right)\)}
\end{problem}

\begin{answer}
B
\end{answer}

\begin{solution}
由参数方程中的绝对值，\(x\geq0\). 对 \(x\) 平方，得
\[
x^2=\left(\cos\frac\theta2+\sin\frac\theta2\right)^2
=1+\sin\theta=2y
\]
所以参数点都在抛物线 \(x^2=2y\) 上，且只取其中满足 \(x\geq0\) 的部分. 当 \(\theta=\pi\) 时，\((x,y)=(1,\frac12)\)，故选 B.
\end{solution}

\begin{problem}
若 \(a\)，\(b\) 是任意实数，且 \(a>b\)，则（\quad）

\choices
    {\(a^2>b^2\)}
    {\(\frac{b}{a}<1\)}
    {\(\lg(a-b)>0\)}
    {\(\left(\frac{1}{2}\right)^a<\left(\frac{1}{2}\right)^b\)}
\end{problem}

\begin{answer}
D
\end{answer}

\begin{solution}
前三个选项均不恒成立：取 \(a=1\)，\(b=-2\)，有 \(a>b\) 但 \(a^2<b^2\)；取 \(a=-1\)，\(b=-2\)，有 \(b/a=2>1\)；取 \(a=1\)，\(b=0\)，有 \(\lg(a-b)=\lg1=0\). 对于选项 D，因为底数 \(\frac12\in(0,1)\)，指数函数 \(y=\left(\frac12\right)^x\) 是减函数，所以由 \(a>b\) 得
\[
\left(\frac12\right)^a<\left(\frac12\right)^b
\]
故选 D.
\end{solution}

\begin{problem}
一动圆与两圆 \(x^2+y^2=1\) 和 \(x^2+y^2-8x+12=0\) 都外切，则动圆圆心轨迹为（\quad）

\choices
    {圆}
    {椭圆}
    {双曲线的一支}
    {抛物线}
\end{problem}

\begin{answer}
C
\end{answer}

\begin{solution}
两已知圆的圆心和半径分别为
\(O_1=(0,0)\)，\(r_1=1;\)，\(O_2=(4,0)\)，\(r_2=2\).
设动圆圆心为 \(P\)，半径为 \(r\). 外切条件给出
\(PO_1=r+1\)，\(PO_2=r+2\).
相减得
\[
PO_2-PO_1=1
\]
这表示点 \(P\) 到两个定点 \(O_1\)，\(O_2\) 的距离之差为定值 \(1\)，且 \(1<O_1O_2=4\)，所以 \(P\) 的轨迹是双曲线的一支，故选 C.
\end{solution}

\begin{problem}
圆柱轴截面的周长 \(l\) 为定值，那么圆柱体积的最大值是（\quad）

\choices
    {\(\left(\frac{l}{6}\right)^3\pi\)}
    {\(\frac{1}{9}\left(\frac{l}{2}\right)^3\pi\)}
    {\(\left(\frac{l}{4}\right)^3\pi\)}
    {\(2\left(\frac{l}{4}\right)^3\pi\)}
\end{problem}

\begin{answer}
A
\end{answer}

\begin{solution}
设圆柱底面半径为 \(r\)，高为 \(h\). 轴截面是长为 \(h\)、宽为 \(2r\) 的矩形，因此
\(2h+4r=l\)，\(h=\frac l2-2r\)，\(0<r<\frac l4\).
体积为
\[
V=\pi r^2h=\pi r^2\left(\frac l2-2r\right)
\]
求导得
\[
V'(r)=\pi r(l-6r)
\]
在 \(0<r<\frac l4\) 内，\(V'(r)\) 在 \(r=\frac l6\) 左侧为正、右侧为负，因此最大值在 \(r=\frac l6\) 时取得. 此时 \(h=\frac l6\)，故
\[
V_{\max}=\pi\left(\frac l6\right)^3
\]
故选 A.
\end{solution}

\begin{problem}
\((\sqrt{x}+1)^4(x-1)^5\) 展开式中 \(x^4\) 的系数为（\quad）

\choices
    {\(-40\)}
    {\(10\)}
    {\(40\)}
    {\(45\)}
\end{problem}

\begin{answer}
D
\end{answer}

\begin{solution}
将 \((\sqrt{x}+1)^4\) 展开，和 \((x-1)^5\) 相乘时，能产生 \(x^4\) 的只有下列三项：
\(x^2\cdot(-10x^2)\)，\(6x\cdot10x^3\)，\(1\cdot(-5x^4)\).
因此 \(x^4\) 的系数为
\[
-10+60-5=45
\]
故选 D.
\end{solution}

\begin{problem}
直角梯形的一个内角为 \(45^\circ\)，下底长为上底长的 \(\frac{3}{2}\)，这个梯形绕下底所在的直线旋转一周所成的旋转体的全面积为 \(\left(5+\sqrt{2}\right)\pi\)，则旋转体的体积为（\quad）

\choices
    {\(2\pi\)}
    {\(\frac{4+\sqrt{2}}{3}\pi\)}
    {\(\frac{5+\sqrt{2}}{3}\pi\)}
    {\(\frac{7}{3}\pi\)}
\end{problem}

\begin{answer}
D
\end{answer}

\begin{solution}
设上底长为 \(a\)，下底长为 \(\frac32a\)，梯形高为 \(h\). 由于一个锐角为 \(45^\circ\)，两底长度之差等于高，所以
\[
h=\frac32a-a=\frac a2
\]
绕下底旋转后，所得旋转体由半径为 \(h\)、高为 \(a\) 的圆柱和半径为 \(h\)、高为 \(\frac a2=h\) 的圆锥组成. 圆柱与圆锥的公共圆面为内部面，外部全面积还包括左端圆面，因此
\[
S=2\pi ah+\pi h\cdot h\sqrt2+\pi h^2
\]
代入 \(h=\frac a2\)，由题设得
\[
(5+\sqrt2)\pi
=\pi a^2\left(1+\frac{\sqrt2+1}{4}\right)
=\frac{5+\sqrt2}{4}\pi a^2
\]
所以 \(a=2\)，\(h=1\). 旋转体体积为
\[
V=\pi h^2a+\frac13\pi h^2\cdot\frac a2
=2\pi+\frac13\pi=\frac73\pi
\]
故选 D.
\end{solution}

\begin{problem}
已知 \(a_1\)，\(a_2\)，\(\cdots\)，\(a_8\) 为各项都大于零的等比数列，公比 \(q\neq 1\)，则（\quad）

\choices
    {\(a_1+a_8>a_4+a_5\)}
    {\(a_1+a_8<a_4+a_5\)}
    {\(a_1+a_8=a_4+a_5\)}
    {\(a_1+a_8\) 和 \(a_4+a_5\) 的大小关系不能由已知条件确定}
\end{problem}

\begin{answer}
A
\end{answer}

\begin{solution}
因各项均为正，公比 \(q>0\). 写出各项得
\[
a_1+a_8-a_4-a_5
=a_1(1+q^7-q^3-q^4)
=a_1(1-q^3)(1-q^4)
\]
当 \(0<q<1\) 时，两个括号均为正；当 \(q>1\) 时，两个括号均为负. 又 \(q\neq1\)，所以两括号的乘积始终为正，从而
\[
a_1+a_8>a_4+a_5
\]
故选 A.
\end{solution}

\begin{problem}
设有如下三个命题：甲：相交两直线 \(l\)，\(m\) 都在平面 \(\alpha\) 内，并且都不在平面 \(\beta\) 内. 乙：\(l\)，\(m\) 之中至少有一条与 \(\beta\) 相交. 丙：\(\alpha\) 与 \(\beta\) 相交. 当甲成立时（\quad）

\choices
    {乙是丙的充分而不必要的条件}
    {乙是丙的必要而不充分的条件}
    {乙是丙的充分且必要的条件}
    {乙既不是丙的充分条件又不是丙的必要条件}
\end{problem}

\begin{answer}
C
\end{answer}

\begin{solution}
设 \(\alpha\cap\beta=a\). 若丙成立，则 \(\alpha\) 与 \(\beta\) 的公共点全在 \(a\) 上. 在平面 \(\alpha\) 内，相交直线 \(l\)，\(m\) 不可能同时都平行于 \(a\)，所以至少有一条直线不平行于 \(a\)，从而与 \(a\subset\beta\) 相交，乙成立.

反过来，若乙成立，例如 \(l\) 与 \(\beta\) 相交，则交点既在 \(l\subset\alpha\) 内，又在 \(\beta\) 内，因此该交点属于 \(\alpha\cap\beta\)，丙成立. 因而在甲成立的条件下，乙与丙等价，故选 C.
\end{solution}

\begin{problem}
将数字 \(1\)，\(2\)，\(3\)，\(4\) 填入标号为 \(1\)，\(2\)，\(3\)，\(4\) 的四个方格里，每格填一个数字，则每个方格的标号与所填的数字均不相同的填法有（\quad）

\choices
    {\(6\) 种}
    {\(9\) 种}
    {\(11\) 种}
    {\(23\) 种}
\end{problem}

\begin{answer}
B
\end{answer}

\begin{solution}
先任意排列 \(1\)，\(2\)，\(3\)，\(4\)，共有 \(4!=24\) 种. 用容斥原理排除至少一个数字与方格标号相同的情况：
\[
4!-\mathrm C_4^1 3!+\mathrm C_4^2 2!-\mathrm C_4^3 1!+\mathrm C_4^4
=24-24+12-4+1=9
\]
所以满足条件的填法有 \(9\) 种，故选 B.
\end{solution}

\section{填空题}

\begin{problem}
\(\sin\left(\arccos\frac{1}{2}+\arccos\frac{1}{3}\right)=\)\fillinblank{}.
\end{problem}

\begin{answer}
\(\frac{\sqrt{3}+2\sqrt{2}}{6}\)
\end{answer}

\begin{solution}
设
\(\alpha=\arccos\frac12\)，\(\beta=\arccos\frac13\).
因为 \(\alpha\)，\(\beta\in[0,\pi]\)，所以
\(\sin\alpha=\sqrt{1-\cos^2\alpha}=\frac{\sqrt3}{2}\)，\(\sin\beta=\sqrt{1-\cos^2\beta}=\frac{2\sqrt2}{3}\).
由正弦的和角公式，
\[
\sin(\alpha+\beta)
=\sin\alpha\cos\beta+\cos\alpha\sin\beta
=\frac{\sqrt3}{2}\cdot\frac13+\frac12\cdot\frac{2\sqrt2}{3}
=\frac{\sqrt3+2\sqrt2}{6}
\]
\end{solution}

\begin{problem}
若双曲线 \(\frac{x^2}{9k^2}-\frac{y^2}{4k^2}=1\) 与圆 \(x^2+y^2=1\) 没有公共点，则实数 \(k\) 的取值范围为\fillinblank{}.
\end{problem}

\begin{answer}
\(\left|k\right|>\frac{1}{3}\)
\end{answer}

\begin{solution}
若双曲线与单位圆有公共点，则在 \(x^2+y^2=1\) 的条件下还应满足
\[
\frac{x^2}{9k^2}-\frac{y^2}{4k^2}=1
\]
将上式乘以 \(36k^2\)，并利用 \(y^2=1-x^2\)，得
\(4x^2-9(1-x^2)=36k^2\)，\(13x^2=36k^2+9\).
因此公共点存在的充要条件是由此得到的 \(x^2\) 不超过 \(1\)，即
\[
\frac{36k^2+9}{13}\leq1
\Longleftrightarrow
36k^2\leq4
\Longleftrightarrow
|k|\leq\frac13
\]
等号时两曲线相切，仍有公共点. 所以没有公共点时
\[
|k|>\frac13
\]
\end{solution}

\begin{problem}
从 \(1\)，\(2\)，\(\cdots\)，\(10\) 这十个数中取出四个数，使它们的和为奇数，共有\fillinblank{} 种取法（用数字作答）.
\end{problem}

\begin{answer}
100
\end{answer}

\begin{solution}
从 \(1\) 到 \(10\) 中有 \(5\) 个奇数和 \(5\) 个偶数. 四个数的和为奇数，当且仅当取出的奇数个数为 \(1\) 或 \(3\). 因此取法数为
\[
\mathrm C_5^1\mathrm C_5^3+\mathrm C_5^3\mathrm C_5^1
=5\cdot10+10\cdot5=100
\]
\end{solution}

\begin{problem}
设 \(f(x)=4^x-2^{x+1}\)，则 \(f^{-1}(0)=\)\fillinblank{}.
\end{problem}

\begin{answer}
1
\end{answer}

\begin{solution}
令 \(t=2^x\)，则 \(t>0\)，且 \(4^x=t^2\). 方程 \(f(x)=0\) 化为
\[
t^2-2t=0
\Longleftrightarrow
t(t-2)=0
\]
由于 \(t>0\)，只能有 \(t=2\)，于是 \(2^x=2\)，解得 \(x=1\). 因而填 \(1\).
\end{solution}

\begin{problem}
建造一个容积为 \(8\mathrm{~m}^3\)，深为 \(2\mathrm{~m}\) 的长方体无盖水池. 如果池底和池壁的造价每平方米分别为 \(120\text{ 元}\) 和 \(80\text{ 元}\)，那么水池的最低总造价为\fillinblank{}\(\text{元}\).
\end{problem}

\begin{answer}
1760
\end{answer}

\begin{solution}
设池底的长、宽分别为 \(x\mathrm{~m}\)，\(y\mathrm{~m}\). 由容积为 \(8\mathrm{~m}^3\)、深为 \(2\mathrm{~m}\)，得
\(2xy=8\)，\(xy=4\).
池底造价为 \(120xy=480\) 元，四面池壁的面积为
\[
2\cdot2x+2\cdot2y=4(x+y)\mathrm{~m}^2
\]
故池壁造价为 \(320(x+y)\) 元，总造价
\[
C=480+320(x+y)
\]
由基本不等式，
\[
x+y\geq2\sqrt{xy}=4
\]
且当 \(x=y=2\) 时取等号. 所以最低总造价为
\[
C_{\min}=480+320\cdot4=1760\text{ 元}
\]
\end{solution}

\begin{problem}
如图，\(ABCD\) 是正方形，\(E\) 是 \(AB\) 的中点，如将 \(\triangle DAE\) 和 \(\triangle CBE\) 分别沿虚线 \(DE\) 和 \(CE\) 折起，使 \(AE\) 与 \(BE\) 重合，记 \(A\) 与 \(B\) 重合的点为 \(P\)，则面 \(PCD\) 与面 \(ECD\) 所成的二面角为\fillinblank{} 度.

\bitmapfigure[max width=0.62\linewidth]{img/1993/jing_xiang_e_yun_qiong_gui_science/q23_fig1.png}
\FloatBarrier
\end{problem}

\begin{answer}
30
\end{answer}

\begin{solution}
设正方形边长为 \(2\)，取 \(CD\) 的中点为 \(M\). 折起后，\(PE=AE=BE=1\)，且
\(PD=DA=2\)，\(PC=CB=2\).
于是 \(PD=PC\)，\(ED=EC\)，从而在等腰三角形 \(PCD\) 和 \(ECD\) 中，\(PM\perp CD\)，\(EM\perp CD\). 所以所求二面角的平面角为 \(\angle PME\).

在三角形 \(PCD\) 中，\(PC=PD=CD=2\)，故它是等边三角形，
\[
PM=\sqrt{2^2-1^2}=\sqrt3
\]
又 \(EM\) 是正方形的高，\(EM=2\)，并且 \(PE=1\). 在三角形 \(PME\) 中，
\[
\cos\angle PME
=\frac{PM^2+EM^2-PE^2}{2\cdot PM\cdot EM}
=\frac{3+4-1}{4\sqrt3}
=\frac{\sqrt3}{2}
\]
因此 \(\angle PME=30^\circ\)，所求二面角为 \(30^\circ\).
\end{solution}

\section{解答题}

\begin{problem}
已知 \(f(x)=\log_a\frac{1+x}{1-x}\)（\(a>0\)，\(a\neq 1\)）.

\begin{enumerate}
\item 求 \(f(x)\) 的定义域；
\item 判断 \(f(x)\) 的奇偶性并予以证明；
\item 求使 \(f(x)>0\) 的 \(x\) 取值范围.
\end{enumerate}
\end{problem}

\begin{solution}
\begin{enumerate}
\item 由对数真数必须为正，得 \(\frac{1+x}{1-x}>0\).
在临界点 \(-1\)，\(1\) 两侧讨论符号，只有 \(-1<x<1\) 时上式成立，所以 \(f(x)\) 的定义域为 \((-1,1)\).

\item 定义域关于原点对称，且对任意 \(x\in(-1,1)\)，有
\[
f(-x)=\log_a\frac{1-x}{1+x}=\log_a\left(\frac{1+x}{1-x}\right)^{-1}=-\log_a\frac{1+x}{1-x}=-f(x)
\]
因此 \(f(x)\) 是奇函数.

\item 在定义域 \((-1,1)\) 内，
\[
\frac{1+x}{1-x}>1\Longleftrightarrow 1+x>1-x\Longleftrightarrow x>0
\]
并且 \(\frac{1+x}{1-x}<1\) 等价于 \(x<0\). 当 \(a>1\) 时，\(\log_a t>0\) 等价于 \(t>1\)，故 \(f(x)>0\) 的解为 \((0,1)\)；当 \(0<a<1\) 时，\(\log_a t>0\) 等价于 \(0<t<1\)，故解为 \((-1,0)\).
\end{enumerate}
\end{solution}

\begin{problem}
已知数列 \(\frac{8\cdot 1}{1^2\cdot 3^2}\)，\(\frac{8\cdot 2}{3^2\cdot 5^2}\)，\(\cdots\)，\(\frac{8n}{(2n-1)^2(2n+1)^2}\)，\(\cdots\). \(S_n\) 为其前 \(n\) 项和. 计算得 \(S_1=\frac{8}{9}\)，\(S_2=\frac{24}{25}\)，\(S_3=\frac{48}{49}\)，\(S_4=\frac{80}{81}\). 观察上述结果，推测出计算 \(S_n\) 的公式，并用数学归纳法加以证明.
\end{problem}

\begin{solution}
由已给出的前四项可猜想
\[
S_n=\frac{4n(n+1)}{(2n+1)^2}=1-\frac{1}{(2n+1)^2}
\]
当 \(n=1\) 时，
\[
S_1=\frac{8\cdot1}{1^2\cdot3^2}=\frac{8}{9}=\frac{4\cdot1\cdot2}{3^2}
\]
所以公式对 \(n=1\) 成立.

假设当 \(n=k\) 时公式成立，即
\[
S_k=\frac{4k(k+1)}{(2k+1)^2}
\]
数列的第 \(k+1\) 项为
\[
a_{k+1}=\frac{8(k+1)}{(2k+1)^2(2k+3)^2}
\]
于是
\[
S_{k+1}=S_k+a_{k+1}
=\frac{4(k+1)\bigl(k(2k+3)^2+2\bigr)}{(2k+1)^2(2k+3)^2}
=\frac{4(k+1)(k+2)}{(2k+3)^2}
\]
最后一步是因为
\[
k(2k+3)^2+2=(k+2)(2k+1)^2
\]
因此，若公式对 \(n=k\) 成立，则对 \(n=k+1\) 也成立. 根据数学归纳法，对任意正整数 \(n\)，都有
\[
\boxed{S_n=\frac{4n(n+1)}{(2n+1)^2}}
\]
\end{solution}

\begin{problem}
已知：平面 \(\alpha\) 与平面 \(\beta\) 相交于直线 \(a\). \(\alpha\)，\(\beta\) 同垂直于平面 \(\gamma\)，又同平行于直线 \(b\). 求证：

\begin{enumerate}
\item \(a\perp\gamma\)；
\item \(b\perp\gamma\).
\end{enumerate}
\end{problem}

\begin{solution}
\begin{enumerate}
\item 取平面 \(\alpha\)，\(\beta\)，\(\gamma\) 的法向量分别为 \(\bs n_\alpha\)，\(\bs n_\beta\)，\(\bs n_\gamma\). 由 \(\alpha\perp\gamma\) 和 \(\beta\perp\gamma\)，有 \(\bs n_\alpha\perp\bs n_\gamma\)，\(\bs n_\beta\perp\bs n_\gamma\)，所以 \(\bs n_\alpha\)，\(\bs n_\beta\) 都平行于平面 \(\gamma\). 直线 \(a=\alpha\cap\beta\) 的方向同时垂直于 \(\bs n_\alpha\) 和 \(\bs n_\beta\)，因 \(\alpha\)，\(\beta\) 相交，两法向量不平行，故 \(a\) 的方向平行于 \(\bs n_\gamma\). 因而 \(a\perp\gamma\).

\item 直线 \(b\) 同时平行于平面 \(\alpha\) 和平面 \(\beta\)，所以 \(b\) 的方向是两平面的公共方向. 两相交平面的公共方向正是交线 \(a\) 的方向，故 \(b\parallel a\). 由第（1）问 \(a\perp\gamma\)，得到 \(b\perp\gamma\).
\end{enumerate}
\end{solution}

\begin{problem}
在面积为 \(1\) 的 \(\triangle PMN\) 中，\(\tan \angle PMN=\frac{1}{2}\)，\(\tan \angle MNP=-2\). 建立适当的坐标系，求以 \(M\)，\(N\) 为焦点且过点 \(P\) 的椭圆方程.
\end{problem}

\begin{solution}
建立直角坐标系，使 \(M=(0,0)\)，\(N=(d,0)\)，并取 \(P\) 在 \(x\) 轴上方. 由 \(\tan\angle PMN=\frac12\)，直线 \(MP\) 的斜率为 \(\frac12\). 由于射线 \(NM\) 的方向为负 \(x\) 轴方向，且 \(\tan\angle MNP=-2\)，射线 \(NP\) 的斜率为 \(2\). 因此 \(P\) 满足
\[
\begin{cases}
y=\dfrac{x}{2},\\
y=2(x-d).
\end{cases}
\]
解得
\[
P=\left(\frac{4d}{3},\frac{2d}{3}\right)
\]
由三角形面积为 \(1\)，有
\[
1=\frac12\cdot MN\cdot\text{高}=\frac12\cdot d\cdot\frac{2d}{3}=\frac{d^2}{3}
\]
所以 \(d=\sqrt3\).

于是
\[
MP=\sqrt{\left(\frac{4d}{3}\right)^2+\left(\frac{2d}{3}\right)^2}=\frac{2\sqrt5}{3}d=\frac{2\sqrt{15}}{3}
\]
\[
NP=\sqrt{\left(\frac{d}{3}\right)^2+\left(\frac{2d}{3}\right)^2}=\frac{\sqrt5}{3}d=\frac{\sqrt{15}}{3}
\]
因此 \(MP+NP=\sqrt{15}\). 以 \(M\)，\(N\) 为焦点的椭圆满足 \(2a=MP+NP\)，故 \(a=\frac{\sqrt{15}}{2}\). 将坐标原点移到 \(MN\) 的中点，则两焦点为 \(\left(-\frac{\sqrt3}{2},0\right)\)，\(\left(\frac{\sqrt3}{2},0\right)\)，从而 \(c=\frac{\sqrt3}{2}\). 于是
\[
b^2=a^2-c^2=\frac{15}{4}-\frac{3}{4}=3
\]
所求椭圆方程为
\[
\boxed{\frac{x^2}{15/4}+\frac{y^2}{3}=1}
\]
\end{solution}

\begin{problem}
设复数 \(z=\cos\theta+\i\sin\theta\)（\(0<\theta<\pi\)），\(\omega=\frac{1-\left(\overline{z}\right)^4}{1+z^4}\)，并且 \(|\omega|=\frac{\sqrt{3}}{3}\)，\(\arg \omega<\frac{\pi}{2}\)，求 \(\theta\).
\end{problem}

\begin{solution}
因为 \(z=\e^{\i\theta}\)，所以 \(\overline z=\e^{-\i\theta}\)，从而
\[
\omega=\frac{1-\e^{-4\i\theta}}{1+\e^{4\i\theta}}
=\frac{2\i\e^{-2\i\theta}\sin2\theta}{2\e^{2\i\theta}\cos2\theta}
=\i\e^{-4\i\theta}\tan2\theta
=\tan2\theta\left(\sin4\theta+\i\cos4\theta\right)
\]
因此
\[
|\omega|=|\tan2\theta|=\frac{\sqrt3}{3}
\]
即 \(\tan2\theta=\pm\frac{\sqrt3}{3}\). 结合 \(0<\theta<\pi\)，得到四个候选值
\(\theta=\frac{\pi}{12}\)，\(\frac{5\pi}{12}\)，\(\frac{7\pi}{12}\)，\(\frac{11\pi}{12}\).
当 \(\theta=\frac{\pi}{12}\) 或 \(\frac{7\pi}{12}\) 时，\(\tan2\theta=\frac{\sqrt3}{3}\)，且 \(4\theta\) 与 \(\frac{\pi}{3}\) 同余，故
\(\omega=\frac12+\frac{\sqrt3}{6}\i\)，\(\arg\omega=\frac{\pi}{6}<\frac{\pi}{2}\).
当 \(\theta=\frac{5\pi}{12}\) 或 \(\frac{11\pi}{12}\) 时，\(\tan2\theta=-\frac{\sqrt3}{3}\)，且 \(4\theta\) 与 \(\frac{5\pi}{3}\) 同余，故
\(\omega=\frac12-\frac{\sqrt3}{6}\i\)，\(\arg\omega=\frac{11\pi}{6}>\frac{\pi}{2}\).
所以
\[
\boxed{\theta=\frac{\pi}{12}\text{ 或 }\theta=\frac{7\pi}{12}}
\]
\end{solution}
